\SourcePage{179}{182}
\section{Time-dependent perturbations}

We now turn to perturbations with an explicit dependence on time. In this case
there can be no question at all of corrections to the energy eigenvalues:
when the Hamiltonian depends on time (as does the perturbed operator
$\hat H=\hat H_0+\hat V(t)$), energy is not conserved, and stationary states
therefore do not exist. The problem is instead to calculate approximate wave
functions from the wave functions of the unperturbed system's stationary
states.

For this purpose we use the analogue of the familiar variation-of-constants
method for solving linear differential equations ($P.\,A.\,M.\,Dirac$, 1926).
Let $\Psi_k^{(0)}$ be the wave functions (including the time factor) of the
unperturbed system's stationary states. An arbitrary solution of the
unperturbed wave equation can then be written as
$\Psi=\sum_k a_k\Psi_k^{(0)}$. We seek a solution of the perturbed equation
\begin{equation}
  i\hbar\frac{\partial\Psi}{\partial t}=(\hat H_0+\hat V)\Psi
  \tag{40.1}
\end{equation}
in the form
\begin{equation}
  \Psi=\sum_k a_k(t)\Psi_k^{(0)},
  \tag{40.2}
\end{equation}
where the expansion coefficients are functions of time. Substitution of
(40.2) into (40.1), with
\[
  i\hbar\frac{\partial\Psi_k^{(0)}}{\partial t}=\hat H_0\Psi_k^{(0)},
\]
gives
\[
  i\hbar\sum_k\Psi_k^{(0)}\frac{da_k}{dt}
  =\sum_k a_k\hat V\Psi_k^{(0)}.
\]
\SourcePage{180}{183}
Multiplying both sides on the left by $\Psi_m^{(0)*}$ and integrating, we find
\begin{equation}
  i\hbar\frac{da_m}{dt}=\sum_k V_{mk}(t)a_k,
  \tag{40.3}
\end{equation}
where
\[
  V_{mk}(t)=\int\Psi_m^{(0)*}\hat V\Psi_k^{(0)}\,dq
  =V_{mk}e^{i\omega_{mk}t},
  \qquad
  \omega_{mk}=\frac{E_m^{(0)}-E_k^{(0)}}{\hbar}
\]
are the perturbation matrix elements including their time factor. If $V$
itself depends explicitly on time, the quantities $V_{mk}$ must of course
also be regarded as functions of time.

Choose the wave function of the $n$th stationary state as the unperturbed
wave function. The corresponding coefficients in (40.2) are
$a_n^{(0)}=1$ and $a_k^{(0)}=0$ for $k\ne n$. To obtain the first
approximation, write $a_k=a_k^{(0)}+a_k^{(1)}$ and put
$a_k=a_k^{(0)}$ on the right-hand side of (40.3), which already contains the
small quantities $V_{mk}$. Thus
\begin{equation}
  i\hbar\frac{da_k^{(1)}}{dt}=V_{kn}(t).
  \tag{40.4}
\end{equation}
To indicate the unperturbed function for which the correction is being
calculated, introduce a second index on $a_k$ and write
\[
  \Psi_n=\sum_k a_{kn}(t)\Psi_k^{(0)}.
\]
The integrated form of (40.4) is accordingly
\begin{equation}
  a_{kn}^{(1)}=-\frac{i}{\hbar}\int V_{kn}(t)\,dt
  =-\frac{i}{\hbar}\int V_{kn}e^{i\omega_{kn}t}\,dt.
  \tag{40.5}
\end{equation}
This determines the wave functions in the first approximation.

Consider more closely the important case of a time-periodic perturbation
\begin{equation}
  \hat V=\hat F e^{-i\omega t}+\hat G e^{i\omega t},
  \tag{40.6}
\end{equation}
where $\hat F$ and $\hat G$ do not depend on time. Hermiticity of $V$ requires
\[
  \hat F e^{-i\omega t}+\hat G e^{i\omega t}
  =\hat F^{+}e^{i\omega t}+\hat G^{+}e^{-i\omega t},
\]
\SourcePage{181}{184}
whence $\hat G=\hat F^{+}$, or
\begin{equation}
  G_{nm}=F_{mn}^{*}.
  \tag{40.7}
\end{equation}
Using this relation, we have
\begin{equation}
  V_{kn}(t)=V_{kn}e^{i\omega_{kn}t}
  =F_{kn}e^{i(\omega_{kn}-\omega)t}
  +F_{nk}^{*}e^{i(\omega_{kn}+\omega)t}.
  \tag{40.8}
\end{equation}
Substitution in (40.5) and integration yield
\begin{equation}
  a_{kn}^{(1)}
  =-\frac{F_{kn}e^{i(\omega_{kn}-\omega)t}}
           {\hbar(\omega_{kn}-\omega)}
   -\frac{F_{nk}^{*}e^{i(\omega_{kn}+\omega)t}}
           {\hbar(\omega_{kn}+\omega)}.
  \tag{40.9}
\end{equation}
These expressions apply provided neither denominator vanishes\footnote{More
precisely, neither may be so small that $a_{kn}^{(1)}$ ceases to be small in
comparison with unity.}; that is, for every $k$ (at fixed $n$),
\begin{equation}
  E_k^{(0)}-E_n^{(0)}\ne\pm\hbar\omega.
  \tag{40.10}
\end{equation}

For a number of applications it is useful to know the matrix elements of an
arbitrary quantity $f$, defined using the perturbed wave functions. To first
order,
\[
  f_{nm}(t)=f_{nm}^{(0)}(t)+f_{nm}^{(1)}(t),
\]
where
\[
  f_{nm}^{(0)}(t)=\int\Psi_n^{(0)*}\hat f\Psi_m^{(0)}\,dq
  =f_{nm}^{(0)}e^{i\omega_{nm}t},
\]
\[
  f_{nm}^{(1)}(t)=\int
  \bigl(\Psi_n^{(0)*}\hat f\Psi_m^{(1)}
       +\Psi_n^{(1)*}\hat f\Psi_m^{(0)}\bigr)\,dq.
\]
Inserting
\[
  \Psi_n^{(1)}=\sum_k a_{kn}^{(1)}\Psi_k^{(0)}
\]
with the coefficients from (40.9), we obtain
\begin{equation}
\begin{aligned}
  f_{nm}^{(1)}(t)=-e^{i\omega_{nm}t}\sum_k\Bigg\{&
  \left[
    \frac{f_{nk}^{(0)}F_{km}}{\hbar(\omega_{km}-\omega)}
    +\frac{f_{km}^{(0)}F_{nk}}{\hbar(\omega_{kn}+\omega)}
  \right]e^{-i\omega t} \\
  {}+&\left[
    \frac{f_{nk}^{(0)}F_{mk}^{*}}{\hbar(\omega_{km}+\omega)}
    +\frac{f_{km}^{(0)}F_{kn}^{*}}{\hbar(\omega_{kn}-\omega)}
  \right]e^{i\omega t}\Bigg\}.
\end{aligned}
  \tag{40.11}
\end{equation}
This formula applies if none of its terms becomes large, hence if all the
frequencies $\omega_{kn}$ and $\omega_{km}$ remain sufficiently far from
$\omega$. At $\omega=0$ it reduces to (38.12).
\SourcePage{182}{185}
The formulas above assume that the unperturbed energy levels form only a
discrete spectrum. They extend immediately to the presence of a continuous
spectrum as well (the states being perturbed still belonging to the discrete
spectrum): one simply adds the corresponding integrals over the continuous
spectrum to the sums over discrete levels. The denominators
$\omega_{kn}\pm\omega$ in (40.9) and (40.11) must then remain nonzero as
$E_k^{(0)}$ ranges through both the discrete and continuous spectra. If, as
is usual, the continuous spectrum lies above all discrete levels, condition
(40.10), for example, must be supplemented by
\begin{equation}
  E_{\min}^{(0)}-E_n^{(0)}>\hbar\omega,
  \tag{40.12}
\end{equation}
where $E_{\min}^{(0)}$ is the energy at the lower edge of the continuous
spectrum.

\ProblemHeading

\textbf{1.} Determine how the $n$th and $m$th solutions of the Schrödinger
equation change under a periodic perturbation of the form (40.6), whose
frequency $\omega$ satisfies
$E_m^{(0)}-E_n^{(0)}=\hbar(\omega+\varepsilon)$, where $\varepsilon$ is
small.

\SolutionHeading The method developed above does not apply, since the
coefficient $a_{mn}^{(1)}$ in (40.9) becomes large. Return to the exact
equations (40.3), with $V_{mk}(t)$ given by (40.8). Evidently, the principal
effect comes from terms on their right-hand sides whose time dependence is
governed by the small frequency $\omega_{mn}-\omega$. Dropping all other
terms gives
\[
  i\hbar\frac{da_m}{dt}=F_{mn}e^{i\varepsilon t}a_n,
  \qquad
  i\hbar\frac{da_n}{dt}=F_{mn}^{*}e^{-i\varepsilon t}a_m.
\]
Put $a_ne^{i\varepsilon t}=b_n$. Then
\[
  i\hbar\dot a_m=F_{mn}b_n,
  \qquad
  i\hbar(\dot b_n-i\varepsilon b_n)=F_{mn}^{*}a_m,
\]
and elimination of $a_m$ gives
\[
  \ddot b_n-i\varepsilon\dot b_n+(1/\hbar^2)|F_{mn}|^2b_n=0.
\]
Two independent solutions may be chosen as
\begin{equation}
  a_n=Ae^{i\alpha_1t},\qquad
  a_m=-A\frac{\hbar\alpha_1}{F_{mn}}e^{i\alpha_2t}
  \tag{1}
\end{equation}
and
\begin{equation}
  a_n=Be^{-i\alpha_2t},\qquad
  a_m=B\frac{\hbar\alpha_2}{F_{mn}^{*}}e^{-i\alpha_1t},
  \tag{2}
\end{equation}
where $A$ and $B$ are constants (to be fixed by normalization), and
\SourcePage{183}{186}
\[
  \alpha_1=-\varepsilon/2+\Omega,\quad
  \alpha_2=\varepsilon/2+\Omega,\quad
  \Omega=\sqrt{\varepsilon^2/4+|\eta|^2},\quad
  \eta=F_{mn}/\hbar.
\]
Thus the perturbation transforms $\Psi_n^{(0)}$ and $\Psi_m^{(0)}$ into
$a_n\Psi_n^{(0)}+a_m\Psi_m^{(0)}$, with $a_n,a_m$ taken from (1) or (2).

Suppose that at $t=0$ the system was in the state $\Psi_m^{(0)}$. At later
times its state is the linear combination of the two solutions just found
which reduces to $\Psi_m^{(0)}$ at $t=0$:
\begin{equation}
  \Psi=e^{i\varepsilon t/2}
  \left(\cos\Omega t-\frac{i\varepsilon}{2\Omega}\sin\Omega t\right)
  \Psi_m^{(0)}
  -\frac{i\eta^{*}}{\Omega}e^{-i\varepsilon t/2}\sin\Omega t\cdot
  \Psi_n^{(0)}.
  \tag{3}
\end{equation}
The squared modulus of the coefficient of $\Psi_n^{(0)}$ is
\begin{equation}
  \frac{|\eta|^2}{2\Omega^2}\,[1-\cos(2\Omega t)].
  \tag{4}
\end{equation}
It is the probability that the system is in $\Psi_n^{(0)}$ at time $t$.
This probability is periodic, with frequency $2\Omega$, and varies from zero
to $|\eta|^2/\Omega^2$.

At exact resonance, $\varepsilon=0$, (4) becomes
\[
  (1/2)[1-\cos(2|\eta|t)].
\]
It varies periodically between zero and unity: the system passes periodically
from $\Psi_m^{(0)}$ to $\Psi_n^{(0)}$.
